\section{Questão 02 -- Pós-Treino com SFT}
\label{sec:q2}

\subsection{Objetivo}

A Questão 02 aplica o ajuste fino supervisionado (SFT) ao \texttt{Qwen2.5-0.5B},
atualizando \emph{todos} os 494 milhões de parâmetros sobre pares de instrução e
resposta derivados do dataset docentesDC. O objetivo é avaliar se o SFT melhora a
capacidade do modelo de responder no estilo instrução/resposta sobre o conteúdo
acadêmico do Departamento de Computação.

\subsection{Geração dos pares de instrução e resposta}

A partir dos registros do docentesDC foram gerados \textbf{3.111 pares} no formato
\texttt{instruction} / \texttt{input} / \texttt{output}, valor que supera com folga o
mínimo de 1.000 exigido pelo enunciado. Cada par associa uma instrução referente ao
material de um professor a um trecho de resposta extraído desse material. Os pares
são então formatados no molde de prompt estilo
Alpaca~\cite{taori2023alpaca}, como ilustra a Listagem~\ref{lst:par-sft}.

\begin{lstlisting}[caption={Exemplo de par formatado para o SFT (estilo Alpaca).},label={lst:par-sft},captionpos=b]
### Instrução:
Explique o que é tratado neste trecho do material do professor
Weslley Emmanuel Martins Lima.

### Resposta:
Ponteiros são tipos especiais de dados que armazenam endereços de
memória. Uma variável do tipo ponteiro deve ser declarada da forma:
tipo *nome_variavel; ...<|endoftext|>
\end{lstlisting}

Os 3.111 pares foram divididos em 2.800 para treino e 311 para avaliação. A
Tabela~\ref{tab:q2-hiper} resume a configuração do experimento.

\begin{table}[h]
  \centering
  \caption{Configuração do SFT (Questão 02).}
  \label{tab:q2-hiper}
  \footnotesize
  \begin{tabularx}{0.92\textwidth}{@{}l X@{}}
    \toprule
    \textbf{Item} & \textbf{Valor} \\
    \midrule
    Pares gerados (docentesDC) & 3.111 (mínimo exigido: 1.000) \\
    Split treino / avaliação & 2.800 / 311 \\
    Parâmetros treinados & 494M (100\%, \textit{full fine-tuning}) \\
    Épocas / passos & 3 / 525 \\
    Formato de prompt & estilo Alpaca (\texttt{\#\#\# Instrução} / \texttt{\#\#\# Resposta}) \\
    \bottomrule
  \end{tabularx}
\end{table}

\subsection{Resultado \textit{baseline} (antes do SFT)}

A avaliação do modelo base sobre o conjunto de 311 exemplos, antes de qualquer
treino, produziu os valores da Tabela~\ref{tab:q2-baseline}. Esses números são reais
e servem de ponto de partida para a comparação pós-treino.

\begin{table}[h]
  \centering
  \caption{\textit{Baseline} da Questão 02: modelo base sobre o conjunto docentesDC, antes do SFT.}
  \label{tab:q2-baseline}
  \begin{tabular}{@{}l c c@{}}
    \toprule
    \textbf{Métrica} & \textbf{Antes (baseline)} & \textbf{Depois (SFT)} \\
    \midrule
    Entropia Cruzada   & 3,3151 & \textit{em andamento} \\
    Perplexidade       & 27,53  & \textit{em andamento} \\
    Acurácia de Tokens & 41,37\% & \textit{em andamento} \\
    \bottomrule
  \end{tabular}
\end{table}

Qualitativamente, antes do SFT o modelo responde de forma vaga e genérica. Para a
pergunta ``Quem é o professor responsável pela disciplina de Tópicos em IA?'', o
modelo base devolve ``\textit{O professor responsável \ldots{} é [Professor X]}'',
evidenciando ausência de conhecimento específico sobre o corpo docente.

\subsection{Estado da execução e pontos de melhoria identificados}

O laço de treinamento do SFT foi iniciado, mas \textbf{interrompido manualmente}
durante a primeira época. A execução tentada em uma GPU local estimava de 7 a 9 horas
para completar os 525 passos, inviável dentro da janela de hardware disponível;
diferentemente da Questão 01, este treino ainda não foi reexecutado na T4 do Kaggle.
Em consequência, as métricas pós-treino, a comparação final e o salvamento do modelo
permanecem pendentes (sinalizados como \textit{em andamento} na
Tabela~\ref{tab:q2-baseline}).

A análise dos artefatos já gerados também revelou três oportunidades de melhoria de
qualidade, a serem aplicadas antes da execução definitiva:

\begin{enumerate}[leftmargin=2.2em]
  \item \textbf{Qualidade dos pares.} Parte das respostas é composta por texto bruto
        (fragmentos de código, trechos em inglês), o que tende a ensinar o modelo a
        regurgitar conteúdo em vez de transferir conhecimento. Recomenda-se
        deduplicar, filtrar idioma e garantir respostas coerentes em português.
  \item \textbf{Máscara de \textit{loss}.} A perda deve ser calculada apenas sobre os
        tokens da resposta (mascarando a instrução com \texttt{-100} até o marcador
        \texttt{\#\#\# Resposta:}), para o modelo aprender a responder e não a
        reproduzir a pergunta.
  \item \textbf{Perguntas de avaliação reais.} Substituir os marcadores genéricos
        (``professor X/Y'') por nomes que de fato constam no dataset (p.\,ex.\
        Raimundo Moura, Erico Leão, André Soares, Vinícius Machado), tornando a
        comparação antes/depois significativa.
\end{enumerate}

O caminho de conclusão desta questão, junto com a Questão 03, é detalhado na
Seção~\ref{sec:consideracoes}.
